\begin{document}
\docTitle{Metrics}

\paragraph{Surprisal}\phantomsection\label{sec:surprisal}\SourceLabel{sec:surprisal}
The surprisal of an outcome \(x\) that we consider that could happen with probability \(p(x)\) is \(I(x)=-\log p(x)\).
It's also called self-information or information content.
A certain outcome has surprisal 0, a rarer outcome a larger one.
Surprisals of independent outcomes add, because the logarithm turns a product of probabilities into a sum.
One bit is the surprisal of an outcome with probability one half, \(-\log_2\tfrac12=1\), the surprise of one fair coin flip.

\paragraph{Average surprisal: entropy and cross entropy}\phantomsection\label{sec:average-surprisal}\SourceLabel{sec:average-surprisal}
Entropy \(H(p)=\mathbb{E}_{x\sim p}[-\log p(x)]\) is the average surprisal of outcomes assuming that we consider them happening the propability from the true distribution \(p\).
The cross-entropy \(H(p,q)=\mathbb{E}_{x\sim p}[-\log q(x)]\) is the average surprisal assuming that the outcomes come from \(p\), but their surprisal is measured with a model \(q\).
Cross entropy is always higher than entropy.

\paragraph{Which distributions have the highest entropy?}\SourceLabel{sec:entropy-extremes}
\agentProposedEdit[created=2026-11-12T11:18:25.340Z,conversation_id=0b6c5b7e-1d2a-5c3b-9f8e-000000000001]{agent-3e82dcc8-8d37-46f2-b463-b50733de259e}{Adds the notes on entropy, entropy over a sequence, Monte Carlo estimation and the minimum and maximum entropy to this document. The minimum is now said without the word "certain".
}{On \(n\) outcomes, the uniform distribution has the highest entropy, \(\log_2 n\) bits, and a single certain outcome has the lowest, 0 bits.
}{On \(n\) outcomes, the uniform distribution has the highest entropy, \(\log_2 n\) bits, and a distribution with probability 1 on one outcome has the lowest, 0 bits.
}With 4 outcomes, the uniform distribution has 2 bits, \((0.7,0.1,0.1,0.1)\) has 1.36 bits, and \((1,0,0,0)\) has 0 bits.
\agentProposedEdit[created=2026-11-12T11:18:25.340Z,conversation_id=0b6c5b7e-1d2a-5c3b-9f8e-000000000001]{agent-44db0af3-836c-4a7a-8aac-ecbfa7282c4b}{Adds the notes on entropy, entropy over a sequence, Monte Carlo estimation and the minimum and maximum entropy to this document. An example: a vocabulary of 128 tokens, a model undecided between two, and the time t to go from 1 to 0.5 when the shrinking is continuous. The maximum is 7 bits, the two-way case is 1 bit = 0.693 nats, and the other 126 tokens add nothing.
}{}{For a model with a vocabulary of 128 tokens, the highest entropy of the next token is \(\log_2 128=7\) bits, or 4.85 nats.
A model undecided between two tokens, each with probability 0.5, has \(H=-(0.5\log_2 0.5+0.5\log_2 0.5)=1\) bit, whatever the size of the vocabulary, because the other 126 tokens add \(0\log 0=0\).
In nats the same uncertainty is \(\ln 2=0.693\): it is the time \(t\) that continuous shrinking \(e^{-t}\) needs to go from 1 to 0.5, while counting halvings \(2^{-t}\) needs exactly \(t=1\).
}The same holds for continuous distributions on a bounded interval: the flatter the curve, the higher the entropy (Figure~\ref{fig:entropy-shapes}).
On the whole real line no highest value exists without a limit, because a distribution can always spread further.
With the variance fixed, the Gaussian has the highest entropy, \(h=\tfrac12\log_2(2\pi e\sigma^2)\) bits, which depends on the variance and not on the mean.
Continuous entropy has no lowest value and can be negative: it goes to \(-\infty\) as \(\sigma\to 0\).
\providecommand{\captionof}[1]{\expandafter\def\csname @captype\endcsname{#1}\caption}
\begin{center}
\includegraphics[width=0.75\linewidth]{agent-figures/metrics-entropy-shapes}
\captionof{figure}{Three distributions on \([0,1]\): a narrow spike has the lowest entropy (\(-3.01\) bits, orange), a wider curve is in the middle (\(-0.50\) bits, green), and the uniform has the highest (0 bits, dashed blue).}\label{fig:entropy-shapes}
\end{center}

\agentProposedEdit[created=2026-11-12T11:18:25.340Z,conversation_id=0b6c5b7e-1d2a-5c3b-9f8e-000000000001]{agent-166c6b96-7147-4a72-95ae-5e3da47b0bc7}{Adds the notes on entropy, entropy over a sequence, Monte Carlo estimation and the minimum and maximum entropy to this document. The tree start -> A -> x and start -> B -> y1..y4, with H = sum p(y)(-log p(y)), drawn as the figure (0.5 * 0.25 = 0.125 for each B path, 2 bits in total). "Monte Carlo estimation: any expectation can be estimated by sampling", E[f(X)] \textasciitilde{} (1/N) sum f(x\textasciicircum{}(i)), converging as N -> infinity by the law of large numbers, with two additions (independent samples from p, error about sigma/sqrt(N)) and the steps and cautions for a language model.
}{}{\paragraph{What is the entropy of a sequence?}\SourceLabel{sec:sequence-entropy}
A model writes a sequence \(y\) token by token, so \(p(y)\) is the product of the probabilities along its path, and its surprisal \(-\log p(y)\) is the sum of the surprisals of its tokens.
The entropy of the sequences is the same average as for one token, now over whole sequences, \(H=\sum_y p(y)\,(-\log p(y))\).
In the tree of Figure~\ref{fig:sequence-tree}, the sequence A~\(x\) has probability \(0.5\cdot 1=0.5\) and surprisal 1 bit, and each of the four sequences B~\(y_i\) has probability \(0.5\cdot 0.25=0.125\) and surprisal 3 bits.
So \(H=0.5\cdot 1+4\cdot 0.125\cdot 3=0.5+1.5=2\) bits, which is also 1 bit for the choice between A and B plus \(0.5\cdot 2\) bits for the choice after B.
\begin{center}
\includegraphics[width=0.75\linewidth]{agent-figures/metrics-sequence-tree}
\captionof{figure}{A model that writes A or B with probability 0.5 each, then \(x\) after A, or one of four tokens after B: the probability of a sequence is the product along its path, and its surprisal is the sum.}\label{fig:sequence-tree}
\end{center}

\paragraph{How is the entropy of a model estimated?}\SourceLabel{sec:monte-carlo-entropy}
With 128 tokens and 100 positions there are \(128^{100}\) sequences, so the sum over all of them cannot be computed.
Entropy is an expectation, \(H=\mathbb{E}_{y\sim p}[-\log p(y)]\), and any expectation can be estimated by sampling, which is called Monte Carlo estimation:
\[
\mathbb{E}[f(X)]\approx\frac{1}{N}\sum_{i=1}^{N} f\bigl(x^{(i)}\bigr),\qquad x^{(i)}\sim p\ \text{independent}.
\]
By the law of large numbers the average converges to the expectation as \(N\to\infty\), and its error shrinks like \(\sigma/\sqrt{N}\), with \(\sigma\) the standard deviation of \(f(X)\).
For a language model, sample \(N\) answers from the model, score each one in one teacher-forced pass, where \(-\log p(y)\) is the sum of the surprisals of its tokens, and average the scores.
In the tree, the four samples A, B, B, A give \((1+3+3+1)/4=2\) bits, and with \(\sigma=1\) bit, 100 samples give about \(\pm 0.1\) bit.
The samples must come from the model being scored, at the same temperature: scoring another model's text, such as a teacher's answers, gives a cross-entropy instead.

}\paragraph{Connection between variance and entropy}
Entropy and variance both measure randomness, but not the same property.
For a yes/no outcome with probability \(p\) both are zero when \(p\) is 0 or 1 and largest at \(p=0.5\), where the entropy is 1 bit and the variance \(p(1-p)\) is 0.25.
Entropy depends only on the probabilities, while variance also depends on how far apart the outcome values lie (Table~\ref{tab:entropy-variance}).
\begin{center}
\captionof{table}{Two fair coins with different payouts: the same entropy, different variances.}\label{tab:entropy-variance}
\begin{tabular}{lrr}
\hline
Outcomes, each with probability 0.5 & Entropy & Variance\\
\hline
0 or 1 & 1 bit & 0.25\\
0 or 10 & 1 bit & 25\\
\hline
\end{tabular}
\end{center}


\paragraph{Does a multi-modal distribution have lower entropy?}\SourceLabel{sec:entropy-modes}
It depends on what is held fixed (Table~\ref{tab:entropy-modes}).
With the same peak width, two peaks beat one by exactly 1 bit, the choice of which peak.
With the same variance, the single Gaussian wins, because two peaks leave the valley nearly empty and squeeze the probability into two narrow strips (Figure~\ref{fig:two-peaks}).
\begin{center}
\captionof{table}{Entropy in bits of Gaussians with \(\sigma=0.1\), with the two peaks at \(\pm 1\).}\label{tab:entropy-modes}
\begin{tabular}{lrr}
\hline
Distribution & Variance & Entropy\\
\hline
One narrow peak & 0.01 & \(-1.27\)\\
Two narrow peaks & 1.01 & \(-0.27\)\\
One Gaussian with the same variance as the two peaks & 1.01 & 2.05\\
\hline
\end{tabular}
\end{center}
\begin{center}
\includegraphics[width=0.9\linewidth]{agent-figures/metrics-two-peaks}
\captionof{figure}{Two distributions with the same variance 1.01, drawn to scale: two narrow peaks (left) and a single Gaussian (right).}\label{fig:two-peaks}
\end{center}

\paragraph{Standard deviation and width.}\SourceLabel{sec:std-width}
For a single Gaussian, \(\sigma\) is the width: it is the distance from the center to the point where the curve changes its bend (Figure~\ref{fig:gaussian-width}).
About 68\% of the probability lies within \(\pm\sigma\) and about 95\% within \(\pm 2\sigma\), and doubling \(\sigma\) adds exactly 1 bit of entropy.
In general, the standard deviation is the typical distance from the overall mean, the square root of the average squared distance.
For the two peaks of Figure~\ref{fig:two-peaks}, the mean sits in the empty valley, so the standard deviation is about 1 although each peak has \(\sigma=0.1\).
The variance adds the spread inside a peak and the squared distance of the peaks from the mean, \(0.01+1=1.01\), so there it measures how far apart the peaks are.
\begin{center}
\includegraphics[width=0.75\linewidth]{agent-figures/metrics-gaussian-width}
\captionof{figure}{A Gaussian with mean \(\mu\): the arrow of length \(2\sigma\) joins the two points where the curve changes its bend, and the shaded area holds about 68\% of the probability.}\label{fig:gaussian-width}
\end{center}

\agentProposedEdit[created=2026-11-09T09:52:27.019Z,conversation_id=0b6c5b7e-1d2a-5c3b-9f8e-000000000002]{agent-c5c91c14-fb72-47d0-91ef-04c46daccb35}{The metrics note had the pieces that connect variance and entropy (the Gaussian entropy formula, doubling sigma adds 1 bit, the yes/no table) in separate paragraphs but never stated the connection: within a family entropy follows from variance, in general it does not, and the Gaussian gives the upper bound.
}{}{\paragraph{How are variance and entropy connected?}\SourceLabel{sec:variance-entropy}
Both measure how random an outcome is, and within one family of shapes the entropy follows from the variance.
For a Gaussian, \(h=\tfrac12\log_2(2\pi e\sigma^2)\) bits, so the entropy grows with the logarithm of the standard deviation, and doubling \(\sigma\) adds 1 bit.
For a yes/no outcome, both are set by the probability \(p\) alone, the entropy \(-p\log_2 p-(1-p)\log_2(1-p)\) and the variance \(p(1-p)\), and both are largest at \(p=0.5\) (Table~\ref{tab:entropy-variance}).
In general there is no fixed relation, because entropy depends only on the probabilities, while variance also depends on how far apart the outcome values lie.
Among all continuous distributions with a given variance the Gaussian has the highest entropy, so \(\tfrac12\log_2(2\pi e\sigma^2)\) is an upper bound: the two narrow peaks with variance 1.01 have \(-0.27\) bits, against 2.05 bits for the Gaussian (Table~\ref{tab:entropy-modes}).

}
\paragraph{What shape does the log-likelihood have?}\SourceLabel{sec:log-likelihood-shape}
The log-likelihood \(\ell(b)\) of yes/no outcomes under a logistic model is a hump over its parameter \(b\), and the loss \(-\ell(b)\) is the same curve flipped into a bowl.
Each outcome adds one curve, \(\log p\) if it was a success and \(\log(1-p)\) if it was a failure, so one outcome is not a hump, and the hump appears once successes and failures are summed.
The hump is a parabola near its peak and straight lines far from it, with slopes given by the numbers of failures and successes.
Its slope is the score and minus its curvature is the Fisher information, both defined in the \SourceLink{tex/source/statistics.tex}{sec:standard-error}{statistics note} (Figure~\ref{fig:score-information}).
% Numbered captions for pictures that stay in place (what the capt-of package does).
\providecommand{\captionof}[1]{\expandafter\def\csname @captype\endcsname{#1}\caption}
\begin{center}
\includegraphics[width=0.85\linewidth]{agent-figures/metrics-score-information}
\captionof{figure}{The log-likelihood of 2 successes and 4 failures, its slope, the score, and minus its curvature, the Fisher information: at the peak \(\hat b=0.69\) the score is zero and the information is 1.33.}\label{fig:score-information}
\end{center}

\subsection*{Perplexity}
Perplexity is defined as the exponentiated average negative log-likelihood per token.
For a target answer \(y = (y_1,\dots,y_T)\) given prompt \(x\), the model assigns a probability to each next token:

\[
\mathrm{NLL}_{\text{tok}}(x,y)
=
-\frac{1}{T}\sum_{t=1}^{T}
\log p_\theta(y_t\mid x,y_{<t}).
\]

The corresponding perplexity is
\[
\mathrm{PPL}(x,y)
=
\exp\left(\mathrm{NLL}_{\text{tok}}(x,y)\right).
\]
This quantity is often abbreviated as \texttt{ppl}.
The lower the better.


\subsection*{Accuracy +/ Confidence Score}
For tasks with binary outcomes, each evaluation is marked as either solved or not solved.
We define
\[
z_i =
\begin{cases}
1, & \text{if evaluation } i \text{ is solved},\\
0, & \text{otherwise}.
\end{cases}
\]
Given \(N\) total evaluations, the accuracy is the empirical success rate:
\[
\hat{p}
=
\frac{1}{N}\sum_{i=1}^{N} z_i.
\]
Equivalently, \(\hat{p}\) is the fraction of solved evaluations:
\[
\hat{p}
=
\frac{\#\text{solved}}{\#\text{total}}.
\]

Since the outcome is binary, accuracy is equal to the mean of the \(0/1\) variables.

For 0-1 accuracy, the clearest choice is usually a 95\% confidence interval, not standard deviation because 0-1 accuracy is a binomial proportion, not a continuous measurement.
To explain this better consider the following example:
Suppose a model answers 80 out of 100 questions correctly.
For binary results, the standard deviation is approximately:
\(\sqrt{0.8(1-0.8)} = 0.40\)
So you might report:
\(80\% \pm 40\%\)
But this is misleading. 
It does not mean the model's true accuracy is somewhere between 40\% and 120\%.
The 40\% describes how much the individual results vary: each answer is either 0 or 1.
In other way That same standard deviation appears whether you tested 10 questions or 10,000 questions, as long as the accuracy stays near 80\%. 
That is why it does not tell you how reliable the estimated accuracy is.
Confidence internval on the other hand 
This is defined as 
\[
\hat p \pm 1.96 \sqrt{\frac{\hat p(1-\hat p)}{n}}
\]

\end{document}
